\documentclass[11pt]{article}

\usepackage[a4paper,margin=30mm]{geometry}
\usepackage{amsmath,amsthm,amsfonts,amssymb,mathtools}

\newcommand{\E}{\mathbb{E}}
\newcommand{\Pp}{\mathbf{P}}
\newcommand{\R}{\mathbf{R}}
\newcommand{\eqdef}{\mathrel{\vcentcolon=}}
\newcommand{\arginf}{\mathop{\mathrm{arg\,inf}}}
\newcommand{\norm}[1]{\lVert#1\rVert}

\newcommand{\rhohu}{\rho}
\newcommand{\Signal}{X}
\newcommand{\Data}{Y}
\newcommand{\IMFs}{S}
\newcommand{\IMFt}{\widetilde S}
\newcommand{\IMFb}{\bar S}
\newcommand{\BiasedSignal}{X^{\mathrm b}}
\newcommand{\IF}{F}
\newcommand{\Mop}{M}

\newtheorem{proposition}{Proposition}
\newtheorem{lemma}[proposition]{Lemma}
\newtheorem{corollary}[proposition]{Corollary}
\theoremstyle{definition}
\newtheorem*{remark}{Remark}
\newtheorem*{gap}{Gap}

\begin{document}
\begin{center}
  {\Large The second smooth-absolute IRMF component\par}
  \vspace{5pt}
  {\large Sampling error, target shift, and biased signal error after one
   subtraction\par}
  \vspace{5pt}
  {\normalsize Draft}
\end{center}

\section{Purpose}

The first-component note bounds, for the first IRMF component, the sampling
error around the population target, the deterministic target shift, and
their sum, the biased signal error.  Those bounds use two facts that the
second component loses: the summands in the score are independent, and the
noise is Gaussian.  The second component is fitted to the residual
\(\Data_u-\IMFt^{(1)}_u\), whose noise part is the original noise minus the
first-step sampling error.  That part is dependent across locations, is not
Gaussian, and need not be centered.

This draft shows that the same three bounds hold at the second step once two
ingredients are replaced.  Bernstein's inequality is replaced by Gaussian
concentration for Lipschitz functions, which needs no independence.  The
Gaussian closed form of the averaged loss is replaced by a Taylor bound that
only uses the first two moments of the residual noise.  The result is a
sampling error of the same order \(\sqrt{z_x/(nh_2)}\) as at step one,
inflated by an explicit amplification constant, plus a carried-over
remainder of order \(z_x/(nh_1)\), plus a target shift that accumulates the
step-one shift.  Nothing contracts; the one-step rate is preserved.

Section~\ref{sec:recursion} sets up the second step and the surrogate data
on which the analysis runs.  Section~\ref{sec:tools} states the two tools.
Sections~\ref{sec:sampling}, \ref{sec:shift}, and~\ref{sec:biased} bound
the three errors.  Section~\ref{sec:gaps} lists what is not proved here,
including the step from two components to \(k\).  Proofs are in
Appendix~\ref{sec:proofs}.  Statements of the first-component note are cited
as \emph{Note 1}.

\section{The second step}\label{sec:recursion}

\paragraph{Step one, recalled.}
Observations \(\Data_u=\Signal_u+\varepsilon_u\) at \(n\) deterministic
locations, with \(\varepsilon_u\sim N(0,\sigma^2)\) independent, and the
loss \(\rho=\rho_H\) with constants
\(L_1=1\), \(L_2=c/H\), \(L_3=c/H^2\), \(c=\sqrt{2/\pi}\), as in Note~1.
The second step needs the first component at every observation location, so
the step-one evaluation set is the design itself and the step-one union
bound uses \(z^{(1)}_x=x+\log(4n)\).  Step one uses normalized weights
\(a^{(1)}_{v,u}\), the window centred at \(u\) and summing to one over
\(v\), with effective sample size \(N_1\) and curvature lower bound
\(\lambda_1\) as in Note~1.  Write
\(\IMFt^{(1)}_u\), \(\IMFs^{(1)}_u\), \(\BiasedSignal{}^{(1)}_u\) for the
observed component, the population target, and the biased signal at step
one, and
\[
  e^{(1)}_u\eqdef\IMFt^{(1)}_u-\IMFs^{(1)}_u
  =\IF_u^{-1}\nabla^{(1)}_u+R^{(1)}_u,
  \qquad
  \Delta^{(1)}_u\eqdef\IMFs^{(1)}_u-\BiasedSignal{}^{(1)}_u
\]
for the step-one sampling error, its expansion, and the step-one shift.
Note~1 provides an event \(\Omega_1(x)\) of probability at least
\(1-e^{-x}\) on which, for every \(u\),
\begin{equation}
  |e^{(1)}_u|\leq\mathcal E^{(1)}_x,\qquad
  |R^{(1)}_u|\leq\mathcal R^{(1)}_x,
  \qquad\text{and, deterministically,}\qquad
  |\Delta^{(1)}_u|\leq\beta^{(1)}.
  \label{eq:step1-recall}
\end{equation}

\paragraph{Step two.}
The second component is fitted to the residual with a smaller bandwidth
\(h_2<h_1\) and normalized weights \(a^{(2)}_{u,t}\), effective sample size
\(N_2\), on an evaluation set \(\mathcal T\) with \(m=|\mathcal T|\) points.
Three fits are defined exactly as in Note~1, with the data, the clean
signal, and the averaged loss replaced by their step-two versions:
\begin{align}
  \Data^{(2)}_u
    &\eqdef\Data_u-\IMFt^{(1)}_u,
    &
  \IMFt^{(2)}_t
    &\eqdef\arginf_z\sum_u\rhohu(\Data^{(2)}_u-z)\,a^{(2)}_{u,t},
  \label{eq:step2-observed}\\
  \Signal^{(2)}_u
    &\eqdef\Signal_u-\BiasedSignal{}^{(1)}_u,
    &
  \BiasedSignal{}^{(2)}_t
    &\eqdef\arginf_z\sum_u\rhohu(\Signal^{(2)}_u-z)\,a^{(2)}_{u,t}.
  \label{eq:step2-biased}
\end{align}
The biased signal \(\BiasedSignal{}^{(2)}\) is the second output of the
algorithm run on the clean signal.  The population target is defined in the
next paragraph, after the residual has been split into its parts.

\paragraph{What the residual contains.}
Subtracting and adding \(\IMFs^{(1)}_u\) and \(\BiasedSignal{}^{(1)}_u\),
\begin{equation}
  \Data^{(2)}_u
  =\underbrace{\Signal^{(2)}_u-\Delta^{(1)}_u}_{\bar\Signal^{(2)}_u}
   +\underbrace{\varepsilon_u-\IF_u^{-1}\nabla^{(1)}_u}_{\bar\varepsilon^{(2)}_u}
   -R^{(1)}_u.
  \label{eq:residual-split}
\end{equation}
The first part \(\bar\Signal^{(2)}_u\) is deterministic: the clean step-two
signal, displaced by the step-one shift.  The second part
\(\bar\varepsilon^{(2)}_u\) is random and centered, since \(\nabla^{(1)}_u\)
is centered, and it is a smooth function of the Gaussian vector
\(\varepsilon\) with an explicitly bounded gradient.  The third part is the
step-one remainder, which is small on \(\Omega_1(x)\) but has no useful
Lipschitz structure.  The analysis therefore runs on the \emph{surrogate
residual}
\begin{equation}
  \bar\Data^{(2)}_u\eqdef\bar\Signal^{(2)}_u+\bar\varepsilon^{(2)}_u
  =\Data^{(2)}_u+R^{(1)}_u,
  \label{eq:surrogate}
\end{equation}
and the step-one remainder is handled afterwards as a data perturbation.
Define the surrogate fit and the population target by
\begin{equation}
  \IMFb^{(2)}_t
    \eqdef\arginf_z\sum_u\rhohu(\bar\Data^{(2)}_u-z)\,a^{(2)}_{u,t},
  \qquad
  \IMFs^{(2)}_t
    \eqdef\arginf_z\sum_u\E\rhohu(\bar\Data^{(2)}_u-z)\,a^{(2)}_{u,t},
  \label{eq:step2-targets}
\end{equation}
where the expectation is over the whole noise vector.  The averaged loss
now depends on \(u\),
\begin{equation}
  \phi_u(r)\eqdef\E\rhohu(r+\bar\varepsilon^{(2)}_u),
  \qquad
  \E\rhohu(\bar\Data^{(2)}_u-z)=\phi_u(\bar\Signal^{(2)}_u-z),
  \label{eq:phi-u}
\end{equation}
which Note~1 already allows.  The population curvature and the centered
score at the target are
\begin{equation}
  \bar\IF^{(2)}_t\eqdef\sum_u\phi_u''(\bar\Signal^{(2)}_u-\IMFs^{(2)}_t)\,a^{(2)}_{u,t},
  \qquad
  \bar\nabla^{(2)}_t\eqdef\sum_u\left\{
    \rhohu'(\bar\Data^{(2)}_u-\IMFs^{(2)}_t)
    -\E\rhohu'(\bar\Data^{(2)}_u-\IMFs^{(2)}_t)\right\}a^{(2)}_{u,t}.
  \label{eq:step2-curvature-score}
\end{equation}

\paragraph{Three errors.}
With \(\Delta^{(2)}_t\eqdef\IMFs^{(2)}_t-\BiasedSignal{}^{(2)}_t\), the
identity of Note~1 holds verbatim,
\begin{equation}
  \IMFt^{(2)}_t-\BiasedSignal{}^{(2)}_t
  =\underbrace{\IMFt^{(2)}_t-\IMFs^{(2)}_t}_{\text{sampling error}}
   +\underbrace{\Delta^{(2)}_t}_{\text{target shift}},
  \label{eq:step2-identity}
\end{equation}
and the sampling error splits further into the surrogate error and the
perturbation,
\begin{equation}
  \IMFt^{(2)}_t-\IMFs^{(2)}_t
  =(\IMFb^{(2)}_t-\IMFs^{(2)}_t)+(\IMFt^{(2)}_t-\IMFb^{(2)}_t).
  \label{eq:sampling-split}
\end{equation}

\section{Two tools}\label{sec:tools}

\begin{lemma}[Sup-norm stability of the local fit]\label{lem:stability}
Let \(a_u\geq0\) sum to one and let
\(m(y)\eqdef\arginf_z\sum_u\rho_H(y_u-z)a_u\) for \(y\in\R^n\).  Then
\(m\) is continuously differentiable,
\[
  \frac{\partial m}{\partial y_u}(y)
  =\frac{\rho_H''(y_u-m(y))\,a_u}{\sum_v\rho_H''(y_v-m(y))\,a_v}\in[0,1],
  \qquad
  \sum_u\frac{\partial m}{\partial y_u}(y)=1,
\]
and for all \(y,y'\in\R^n\),
\begin{equation}
  |m(y')-m(y)|\leq\max_{u:\,a_u>0}|y'_u-y_u|.
  \label{eq:stability}
\end{equation}
\end{lemma}

The fit is a monotone, weighted-average-like functional of the data.  The
lemma is used three times: to pass from the surrogate fit to the observed
fit, to pass from the displaced clean signal to the clean signal, and to
bound the gradient of the step-one component.

\begin{lemma}[Gaussian concentration for Lipschitz functions]\label{lem:tis}
Let \(\varepsilon\sim N(0,\sigma^2I_n)\) and let \(G:\R^n\to\R\) be
differentiable with \(\sup_\varepsilon\norm{\nabla G(\varepsilon)}_2\leq L\).
Then for every \(s>0\),
\[
  \Pp\bigl(|G(\varepsilon)-\E G(\varepsilon)|>s\bigr)
  \leq2\exp\!\left(-\frac{s^2}{2\sigma^2L^2}\right).
\]
\end{lemma}

This is the Tsirelson--Ibragimov--Sudakov inequality; see for instance
Boucheron, Lugosi, and Massart, \emph{Concentration Inequalities},
Theorem~5.6.  It replaces Bernstein's inequality.  Independence of the
summands is never used; only the Gaussian law of the original noise and a
gradient bound.

\section{Sampling error at step two}\label{sec:sampling}

\paragraph{The step-one component as a function of the noise.}
Differentiating \(\nabla^{(1)}_u\) with respect to \(\varepsilon_v\) gives
\[
  \frac{\partial\bar\varepsilon^{(2)}_u}{\partial\varepsilon_v}
  =\delta_{uv}-\Mop_{u,v},
  \qquad
  \Mop_{u,v}\eqdef
  \frac{\rho_H''(\Data_v-\IMFs^{(1)}_u)\,a^{(1)}_{v,u}}{\bar\IF^{(1)}_u}
  \in\left[0,\frac{L_2}{\lambda_1}a^{(1)}_{v,u}\right].
\]
The matrix \(\Mop\) is the robust equivalent kernel of step one: empirical
second derivatives in the numerator, population curvature in the
denominator.  It depends on the noise, but its entries are bounded
deterministically.  Let
\begin{equation}
  \bar a_1\eqdef\max_v\sum_u a^{(1)}_{v,u},
  \qquad
  \gamma_1\eqdef\frac{L_2}{\lambda_1},
  \qquad
  \Gamma_1\eqdef1+\gamma_1\sqrt{\bar a_1}.
  \label{eq:amplification}
\end{equation}
Here \(\bar a_1\) is the largest total weight that a single observation
receives across all step-one windows; it equals one for a regular periodic
design with constant window mass.  Row sums of \(\Mop\) are at most
\(\gamma_1\), column sums at most \(\gamma_1\bar a_1\), so
\(\norm{\Mop}_{\mathrm{op}}\leq\gamma_1\sqrt{\bar a_1}\) and
\(\norm{I-\Mop}_{\mathrm{op}}\leq\Gamma_1\).  Since
\(\lambda_1\leq c/B\leq c/H=L_2\), always \(\gamma_1\geq1\).

\paragraph{Constants.}
Fix a curvature lower bound \(0<\lambda_2\leq\min_t\bar\IF^{(2)}_t\);
Lemma~\ref{lem:curvature} below gives an explicit choice in the small-noise
regime.  For \(x\geq0\) and \(z_x\eqdef x+\log(4m)\), set
\begin{equation}
  s^{(2)}_1\eqdef\sigma L_2\Gamma_1\sqrt{\frac{2z_x}{N_2}},
  \qquad
  s^{(2)}_2\eqdef\sigma L_3\Gamma_1\sqrt{\frac{2z_x}{N_2}},
  \qquad
  d_2\eqdef\frac{s^{(2)}_2}{\lambda_2}.
  \label{eq:step2-constants}
\end{equation}
Compared with the step-one constants of Note~1, the Bernstein linear term
in \(\alpha z_x\) is gone, the variance factor \(\sigma L_2\) is the
Lipschitz one, and the new factor \(\Gamma_1\) pays for the dependence
introduced by the subtraction.

\begin{lemma}[Step-two score and curvature bounds]\label{lem:step2-local}
There is an event \(\Omega_2(x)\) of probability at least \(1-e^{-x}\) on
which, simultaneously for every \(t\in\mathcal T\),
\begin{equation}
  |\bar\nabla^{(2)}_t|\leq s^{(2)}_1,
  \qquad
  \Bigl|\sum_u\rho_H''(\bar\Data^{(2)}_u-\IMFs^{(2)}_t)\,a^{(2)}_{u,t}
        -\bar\IF^{(2)}_t\Bigr|\leq s^{(2)}_2.
  \label{eq:step2-lemma}
\end{equation}
Deterministically, the third-derivative bounds of Note~1, Lemma~4.1(ii)
and~(iii), hold with \(C_3=L_3/\lambda_2\).
\end{lemma}

\begin{proposition}[Step-two sampling error]\label{prop:step2-sampling}
Suppose
\begin{equation}
  d_2<1,\qquad
  \frac{s^{(2)}_1L_3}{\lambda_2^2(1-d_2)^2}<\frac49.
  \label{eq:step2-smallness}
\end{equation}
Then on \(\Omega_2(x)\), simultaneously for all \(t\), the surrogate fit is
unique and
\begin{align}
  |\IMFb^{(2)}_t-\IMFs^{(2)}_t|
    &\leq\mathcal E^{(2)}_x\eqdef\frac{3s^{(2)}_1}{2\lambda_2(1-d_2)},
  \notag\\
  |\IMFb^{(2)}_t-\IMFs^{(2)}_t-(\bar\IF^{(2)}_t)^{-1}\bar\nabla^{(2)}_t|
    &\leq\mathcal R^{(2)}_x
    \eqdef\frac{9L_3(s^{(2)}_1)^2}{8\lambda_2^3(1-d_2)^3}
          +\frac{s^{(2)}_1s^{(2)}_2}{\lambda_2^2(1-d_2)}.
  \label{eq:step2-surrogate-bounds}
\end{align}
On \(\Omega_1(x)\cap\Omega_2(x)\), which has probability at least
\(1-2e^{-x}\), also
\begin{equation}
  |\IMFt^{(2)}_t-\IMFs^{(2)}_t|
    \leq\mathcal E^{(2)}_x+\mathcal R^{(1)}_x.
  \label{eq:step2-sampling-bound}
\end{equation}
\end{proposition}

The expansion and the remainder bound are exactly those of Note~1,
Proposition~4.2, applied with
\(C_1=s^{(2)}_1/\sqrt{\lambda_2}\), \(C_2=s^{(2)}_2/\sqrt{\lambda_2}\),
\(C_3=L_3/\lambda_2\), \(D_{\min}\geq\sqrt{\lambda_2}\), and
\(\delta\leq d_2\).  The only new probabilistic input is
Lemma~\ref{lem:step2-local}.  The term \(\mathcal R^{(1)}_x\) is the price
of the surrogate: the observed residual differs from the surrogate by the
step-one remainder, and Lemma~\ref{lem:stability} passes that difference
through the step-two fit without amplification.

\section{Target shift at step two}\label{sec:shift}

The step-two shift has two sources.  The population target averages the
loss over the step-two noise \(\bar\varepsilon^{(2)}\), which moves the
minimizer as at step one.  In addition, the deterministic part of the
residual is \(\Signal^{(2)}-\Delta^{(1)}\) rather than \(\Signal^{(2)}\):
the step-one shift has become part of the step-two signal.  Introduce the
clean fit of the displaced signal,
\begin{equation}
  \IMFs^{(2,0)}_t\eqdef\arginf_z\sum_u\rho_H(\bar\Signal^{(2)}_u-z)\,a^{(2)}_{u,t},
  \qquad
  \Delta^{(2)}_t
  =\bigl(\IMFs^{(2)}_t-\IMFs^{(2,0)}_t\bigr)
   +\bigl(\IMFs^{(2,0)}_t-\BiasedSignal{}^{(2)}_t\bigr).
  \label{eq:shift-split}
\end{equation}
The second difference is a clean fit applied to two signals that differ by
at most \(\beta^{(1)}\), so Lemma~\ref{lem:stability} bounds it by
\(\beta^{(1)}\).  The first difference is the noise-averaging shift, and it
needs the first two moments of \(\bar\varepsilon^{(2)}\) and a curvature
bound for the averaged loss.

\paragraph{Moments of the step-two noise.}
Since \(\E\varepsilon_u=0\) and \(\E\nabla^{(1)}_u=0\),
\begin{equation}
  \E\bar\varepsilon^{(2)}_u=0,
  \qquad
  \E(\bar\varepsilon^{(2)}_u)^2\leq\nu^2,
  \qquad
  \nu\eqdef\sigma\left(1+\frac{\gamma_1}{\sqrt{N_1}}\right).
  \label{eq:step2-moments}
\end{equation}
The mean is exactly zero because the surrogate drops the step-one
remainder; the observed residual noise \(\bar\varepsilon^{(2)}-R^{(1)}\)
has mean of order \(\mathcal R^{(1)}_x\).  That bias is accounted for in
\eqref{eq:step2-sampling-bound}, not here.

\paragraph{Curvature of the averaged loss.}
Let \(A^{(2)}_t\) be the width of the active range of \(\bar\Signal^{(2)}\)
in the step-two window at \(t\), and let \(L_4\eqdef c/H^3\), which bounds
\(|\rho_H''''|\).  For \(r>0\), put
\begin{equation}
  \kappa_t(r)\eqdef\frac cH\,e^{-(A^{(2)}_t+r)^2/(2H^2)}-\frac{L_4\nu^2}{2}.
  \label{eq:kappa2}
\end{equation}

\begin{lemma}[Step-two curvature and noise-averaging shift]\label{lem:curvature}
Fix \(t\) and \(r>0\) with \(\kappa_t(r)>0\) and
\(L_3\nu^2<2r\kappa_t(r)\).  Then
\begin{equation}
  |\IMFs^{(2)}_t-\IMFs^{(2,0)}_t|\leq\frac{L_3\nu^2}{2\kappa_t(r)}<r,
  \qquad
  \bar\IF^{(2)}_t\geq\kappa_t(r).
  \label{eq:curvature-shift}
\end{equation}
In particular \(\lambda_2=\min_t\kappa_t(r)\) is an admissible curvature
lower bound in Section~\ref{sec:sampling}.
\end{lemma}

The condition is a small-noise condition: it asks that \(\nu\), which is of
order \(\sigma\), be small compared with \(H\) after accounting for the
signal range.  At step one the Gaussian closed form gave a curvature bound
for every \(\sigma\); here the closed form is unavailable, and the Taylor
bound only pays off when \(L_4\nu^2\) is below the clean curvature.

\begin{proposition}[Step-two target shift]\label{prop:step2-shift}
Under the conditions of Lemma~\ref{lem:curvature}, deterministically,
\begin{equation}
  |\Delta^{(2)}_t|\leq\beta^{(2)}_t
  \eqdef\beta^{(1)}+\frac{L_3\nu^2}{2\kappa_t(r)}
  \qquad\text{for every }t.
  \label{eq:step2-shift-bound}
\end{equation}
\end{proposition}

Both terms are of order \(\sigma^2\) for small \(\sigma\), and neither
decreases with \(n\).  The first term is inherited: the step-one shift
passes through the step-two filter with constant at most one.  The second is
new and has the same form as the step-one shift bound of Note~1,
Proposition~4.5, with \(\sigma^2\) replaced by \(\nu^2\) and the Gaussian
curvature bound replaced by \(\kappa_t(r)\).

\section{Biased signal error at step two}\label{sec:biased}

\begin{corollary}[Combined step-two bounds]\label{cor:step2-biased}
Under the conditions of Propositions~\ref{prop:step2-sampling}
and~\ref{prop:step2-shift}, with probability at least \(1-2e^{-x}\),
simultaneously for all \(t\),
\begin{align}
  |\IMFt^{(2)}_t-\IMFs^{(2)}_t|
    &\leq\mathcal E^{(2)}_x+\mathcal R^{(1)}_x,
    \label{eq:final-step2-population}\\
  |\IMFt^{(2)}_t-\BiasedSignal{}^{(2)}_t|
    &\leq\beta^{(2)}_t+\mathcal E^{(2)}_x+\mathcal R^{(1)}_x.
    \label{eq:final-step2-biased}
\end{align}
\end{corollary}

\paragraph{Rates.}
Suppose the kernel-mass conditions of Note~1 hold at both bandwidths with
the same constant \(C_K\), so \(N_j\geq nh_j/C_K\), and suppose
\(d_2\leq1/2\).  Then
\begin{equation}
  \mathcal E^{(2)}_x
    \leq\frac{3\sigma L_2\Gamma_1}{\lambda_2}\sqrt{\frac{2C_Kz_x}{nh_2}},
  \qquad
  \mathcal R^{(1)}_x
    \leq\left\{\frac{9L_3A_1^2}{\lambda_1^3}
          +\frac{2A_1A_2}{\lambda_1^2}\right\}\frac{C_Kz^{(1)}_x}{nh_1},
  \label{eq:step2-rates}
\end{equation}
with \(A_1,A_2\) the step-one rate constants of Note~1.  The step-two
sampling error therefore has the one-step order \(\sqrt{z_x/(nh_2)}\),
multiplied by \(\Gamma_1\), plus a carried-over term of the smaller order
\(z_x/(nh_1)\).  Since \(h_2=h_1/a\), the leading term is larger than the
step-one sampling error by the factor \(\sqrt a\,\Gamma_1\lambda_1/\lambda_2\)
up to the Bernstein linear term.  The target shift is of order \(\sigma^2\)
and does not decrease.

\paragraph{What this says about contraction.}
Nothing contracts from step one to step two.  The sampling error keeps its
rate in the local sample size and is amplified, not reduced, by the
subtraction.  The target shift accumulates.  This matches the linear audit
in the repository, where the step-\(k\) component error has standard
deviation of order \((nh_k)^{-1/2}\) with a stage-dependent constant.  The
constant \(\Gamma_1\) here is an operator-norm bound and is therefore
conservative; the linear audit shows that the exact constant can
\emph{decrease} after the first stage because later components are
band-pass.  Recovering that cancellation would require tracking the Fourier
structure of \(\Mop\), not just its norm.

\section{Gaps}\label{sec:gaps}

\begin{gap}[Curvature bound for large noise]
Lemma~\ref{lem:curvature} needs \(L_4\nu^2\) below the clean curvature,
roughly \(\sigma\lesssim H\).  Step one has no such restriction.  A
truncation argument, bounding \(\E\rho_H''(r+\bar\varepsilon)\) below by
\(\rho_H''\) at \(|r|+\tau\) times \(\Pp(|\bar\varepsilon|\leq\tau)\), gives
a positive bound for every \(\sigma\) but needs a tail bound for
\(\bar\varepsilon^{(2)}\) beyond the second moment, and a separate argument
that \(\IMFs^{(2)}_t\) stays near the active range.  Both are available in
principle from Lemma~\ref{lem:tis} applied to \(\bar\varepsilon^{(2)}_u\),
which is itself Lipschitz in \(\varepsilon\) with constant \(\Gamma_1\),
but are not written out.
\end{gap}

\begin{gap}[Sharpness of the amplification constant]
\(\Gamma_1\) bounds \(\norm{I-\Mop}_{\mathrm{op}}\) through row and column
sums.  It is at least \(2\) and ignores the cancellation between
\(\varepsilon_u\) and its local average.  For the linear filter on a
periodic grid the exact factor is the Fourier multiplier \(1-w_1(\omega)\),
which is small at low frequencies where the step-two kernel concentrates
its energy.  A sharper constant would need a spectral description of the
robust equivalent kernel, which is data dependent.
\end{gap}

\begin{gap}[Choice of population target]
The population target \(\IMFs^{(2)}_t\) is defined through the surrogate
residual, which drops the step-one remainder.  The alternative, minimizing
\(\sum_u\E\rho_H(\Data^{(2)}_u-z)a^{(2)}_{u,t}\) with the observed
residual, differs from \(\IMFs^{(2)}_t\) by at most
\(\E\max_u|R^{(1)}_u|\) after an \(L^1\) version of
Lemma~\ref{lem:stability}, which is of order \(\mathcal R^{(1)}_x\) plus a
tail contribution.  In simulations, the surrogate target is the one to
compute: the clean step-two filter applied to
\(\Signal-\IMFs^{(1)}\), with the loss averaged over the law of
\(\bar\varepsilon^{(2)}\).
\end{gap}

\begin{gap}[From two steps to \(k\)]
The argument iterates formally.  At step \(k\), the surrogate noise is
\(\bar\varepsilon^{(k)}=(I-\Mop^{(k-1)})\cdots(I-\Mop^{(1)})\varepsilon\)
in the differential sense, so the Lipschitz constant becomes
\(\Gamma_{k-1}\leq\prod_{j<k}(1+\gamma_j\sqrt{\bar a_j})\), the carried
remainder becomes \(\sum_{j<k}\mathcal R^{(j)}_x\), and the shift becomes
\(\sum_{j\leq k}\) of noise-averaging terms.  Since \(\gamma_j\geq1\), the
operator-norm product grows geometrically in \(k\), so this route gives no
uniformity in \(k\).  The linear audit indicates the true constants are
bounded, which again points to the spectral structure of the equivalent
kernels.  Also, \(\nu_k\) grows with \(\Gamma_{k-1}\), which tightens the
small-noise condition at each step.
\end{gap}

\begin{gap}[Inherited from Note 1]
Uniformity over the continuum and adaptive choices of \(h\) or \(H\) remain
open.  The oracle correction at step two uses the unknown step-two target,
curvature, and noise law.
\end{gap}

\appendix
\section{Proofs}\label{sec:proofs}

\begin{proof}[Proof of Lemma~\ref{lem:stability}]
Since \(\rho_H''>0\), the objective is strictly convex and \(m(y)\) is the
unique root of \(g(z,y)\eqdef\sum_u\rho_H'(y_u-z)a_u\).  As
\(\partial_zg=-\sum_u\rho_H''(y_u-z)a_u<0\), the implicit function theorem
gives \(m\in C^1\) with
\[
  \frac{\partial m}{\partial y_u}
  =-\frac{\partial_{y_u}g}{\partial_zg}
  =\frac{\rho_H''(y_u-m)a_u}{\sum_v\rho_H''(y_v-m)a_v}.
\]
These are nonnegative and sum to one.  Along the segment
\(y(s)=y+s(y'-y)\), \(s\in[0,1]\),
\[
  \frac{d}{ds}m(y(s))
  =\sum_u\frac{\partial m}{\partial y_u}(y(s))\,(y'_u-y_u)
  \in\Bigl[-\max_{u:a_u>0}|y'_u-y_u|,\ \max_{u:a_u>0}|y'_u-y_u|\Bigr],
\]
and integrating over \(s\) proves \eqref{eq:stability}.
\end{proof}

\begin{proof}[Proof of Lemma~\ref{lem:step2-local}]
Fix \(t\).  The score functional
\(G_t(\varepsilon)\eqdef\sum_u\rho_H'(\bar\Data^{(2)}_u-\IMFs^{(2)}_t)\,a^{(2)}_{u,t}\)
is smooth in \(\varepsilon\), and \(\bar\nabla^{(2)}_t=G_t-\E G_t\).  By
the chain rule and \eqref{eq:residual-split},
\[
  \frac{\partial G_t}{\partial\varepsilon_v}
  =\sum_u a^{(2)}_{u,t}\,\rho_H''(\bar\Data^{(2)}_u-\IMFs^{(2)}_t)
    \,(\delta_{uv}-\Mop_{u,v})
  =\bigl[c^{\top}(I-\Mop)\bigr]_v,
  \qquad
  0\leq c_u\leq L_2a^{(2)}_{u,t}.
\]
Hence \(\norm{\nabla G_t}_2\leq\norm c_2\norm{I-\Mop}_{\mathrm{op}}\).
Now \(\norm c_2\leq L_2\bigl(\sum_u(a^{(2)}_{u,t})^2\bigr)^{1/2}\leq L_2N_2^{-1/2}\),
and, since \(0\leq\Mop_{u,v}\leq\gamma_1a^{(1)}_{v,u}\), the row sums of
\(\Mop\) are at most \(\gamma_1\) and the column sums at most
\(\gamma_1\bar a_1\).  By the Schur bound
\(\norm{\Mop}_{\mathrm{op}}\leq\sqrt{\norm{\Mop}_1\norm{\Mop}_\infty}\leq\gamma_1\sqrt{\bar a_1}\),
so \(\norm{I-\Mop}_{\mathrm{op}}\leq\Gamma_1\).  Lemma~\ref{lem:tis} with
\(L=L_2\Gamma_1N_2^{-1/2}\) and \(s=\sigma L\sqrt{2z}\) gives
\(\Pp(|\bar\nabla^{(2)}_t|>\sigma L_2\Gamma_1\sqrt{2z/N_2})\leq2e^{-z}\).

The curvature functional
\(H_t(\varepsilon)\eqdef\sum_u\rho_H''(\bar\Data^{(2)}_u-\IMFs^{(2)}_t)\,a^{(2)}_{u,t}\)
has \(\E H_t=\bar\IF^{(2)}_t\) by \eqref{eq:phi-u}, and the same
computation with \(\rho_H'''\) in place of \(\rho_H''\) gives
\(\norm{\nabla H_t}_2\leq L_3\Gamma_1N_2^{-1/2}\), hence
\(\Pp(|H_t-\bar\IF^{(2)}_t|>\sigma L_3\Gamma_1\sqrt{2z/N_2})\leq2e^{-z}\).
Take \(z=z_x\) and the union over the two functionals and the \(m\)
evaluation points; the exceptional probability is at most
\(4me^{-z_x}=e^{-x}\).

For the deterministic bounds, \(|\rho_H'''|\leq L_3\) and the step-two
weights sum to one give the empirical third-derivative bound with
\(C_3=L_3/\lambda_2\) in curvature units.  Differentiation under the
expectation in \eqref{eq:phi-u} is justified by the bounded derivatives, so
\(|\phi_u'''|\leq L_3\) and the population-curvature Lipschitz bound of
Note~1, Lemma~4.1(iii), follows by the mean value theorem exactly as there.
\end{proof}

\begin{proof}[Proof of Proposition~\ref{prop:step2-sampling}]
The population objective \(\sum_u\phi_u(\bar\Signal^{(2)}_u-z)a^{(2)}_{u,t}\)
is finite and coercive by the Lipschitz bound \(|\rho_H'|\leq1\) and
\(\E|\bar\Data^{(2)}_u|<\infty\), convex, and has curvature
\(\bar\IF^{(2)}_t\geq\lambda_2>0\) at its minimizer, which is therefore
unique.  The deterministic part of the proof of Note~1,
Proposition~4.2, uses only the three bounds of Note~1, Lemma~4.1, on a
specified event.  On \(\Omega_2(x)\) these hold with
\(C_1=s^{(2)}_1/\sqrt{\lambda_2}\), \(C_2=s^{(2)}_2/\sqrt{\lambda_2}\),
\(C_3=L_3/\lambda_2\), \(D_t\geq\sqrt{\lambda_2}\), and \(\delta\leq d_2\),
by Lemma~\ref{lem:step2-local}.  Condition \eqref{eq:step2-smallness} is
the corresponding smallness condition, and substitution into the
localization and remainder bounds of Note~1 gives
\eqref{eq:step2-surrogate-bounds}, exactly as in the proof of Note~1,
Corollary~4.3.

For \eqref{eq:step2-sampling-bound}, by \eqref{eq:surrogate} the observed
and surrogate residuals differ by \(R^{(1)}_u\).  Lemma~\ref{lem:stability}
applied to the step-two fit gives
\(|\IMFt^{(2)}_t-\IMFb^{(2)}_t|\leq\max_u|R^{(1)}_u|\leq\mathcal R^{(1)}_x\)
on \(\Omega_1(x)\).  Combine with \eqref{eq:sampling-split}.
\end{proof}

\begin{proof}[Proof of \eqref{eq:step2-moments}]
Centering is immediate.  For the second moment, Minkowski's inequality
gives
\(\bigl(\E(\bar\varepsilon^{(2)}_u)^2\bigr)^{1/2}
 \leq\sigma+(\bar\IF^{(1)}_u)^{-1}\bigl(\E(\bar\nabla^{(1)}_u)^2\bigr)^{1/2}\),
where \(\bar\nabla^{(1)}_u\) is the normalized step-one score.  Its
summands are independent and centered, with variances at most
\(\sigma^2L_2^2(a^{(1)}_{v,u})^2\) by the Lipschitz variance bound in the
proof of Note~1, Corollary~4.3.  Hence
\(\E(\bar\nabla^{(1)}_u)^2\leq\sigma^2L_2^2/N_1\), and
\(\bar\IF^{(1)}_u\geq\lambda_1\) gives the claim with
\(\nu=\sigma(1+\gamma_1N_1^{-1/2})\).
\end{proof}

\begin{proof}[Proof of Lemma~\ref{lem:curvature}]
Write \(\psi=\rho_H'\) and let
\(\bar G(z)\eqdef\sum_u a^{(2)}_{u,t}\E\psi(\bar\Signal^{(2)}_u+\bar\varepsilon^{(2)}_u-z)\)
be the averaged score, so that \(\bar G(\IMFs^{(2)}_t)=0\) and
\(-\bar G'=\sum_ua^{(2)}_{u,t}\E\rho_H''(\cdot)\).  Taylor's theorem with
\(|\psi''|\leq L_3\) and \(|\rho_H''''|\leq L_4\), followed by expectation
with \(\E\bar\varepsilon^{(2)}_u=0\) and \(\E(\bar\varepsilon^{(2)}_u)^2\leq\nu^2\),
gives for every \(r_0\in\R\)
\[
  |\E\psi(r_0+\bar\varepsilon^{(2)}_u)-\psi(r_0)|\leq\frac{L_3\nu^2}{2},
  \qquad
  \E\rho_H''(r_0+\bar\varepsilon^{(2)}_u)\geq\rho_H''(r_0)-\frac{L_4\nu^2}{2}.
\]
Let \(m\eqdef\IMFs^{(2,0)}_t\).  Since the clean score vanishes at \(m\)
and the weights sum to one, the first bound gives
\(|\bar G(m)|\leq L_3\nu^2/2\).  For \(|z-m|\leq r\) and active \(u\),
\(|\bar\Signal^{(2)}_u-z|\leq A^{(2)}_t+r\), because \(m\) lies in the
active range; the second bound then gives \(-\bar G'(z)\geq\kappa_t(r)\).
Hence \(\bar G(m+r)\leq L_3\nu^2/2-r\kappa_t(r)<0\) and
\(\bar G(m-r)\geq-L_3\nu^2/2+r\kappa_t(r)>0\).  The strictly decreasing
function \(\bar G\) therefore has its unique zero \(\IMFs^{(2)}_t\) in
\((m-r,m+r)\), and the mean value theorem on the interval between \(m\) and
\(\IMFs^{(2)}_t\) gives
\(|\IMFs^{(2)}_t-m|\leq|\bar G(m)|/\kappa_t(r)\leq L_3\nu^2/(2\kappa_t(r))\).
Finally \(\bar\IF^{(2)}_t=-\bar G'(\IMFs^{(2)}_t)\geq\kappa_t(r)\).
\end{proof}

\begin{proof}[Proof of Proposition~\ref{prop:step2-shift}]
Use the split \eqref{eq:shift-split}.  The first difference is bounded by
Lemma~\ref{lem:curvature}.  For the second, \(\IMFs^{(2,0)}_t\) and
\(\BiasedSignal{}^{(2)}_t\) are the same clean fit applied to
\(\bar\Signal^{(2)}=\Signal^{(2)}-\Delta^{(1)}\) and to \(\Signal^{(2)}\),
and Lemma~\ref{lem:stability} with \eqref{eq:step1-recall} bounds their
difference by \(\max_u|\Delta^{(1)}_u|\leq\beta^{(1)}\).
\end{proof}

\begin{proof}[Proof of Corollary~\ref{cor:step2-biased} and the rates]
Add the deterministic bound of Proposition~\ref{prop:step2-shift} to
\eqref{eq:step2-sampling-bound} through the identity
\eqref{eq:step2-identity}; the failure probability is unchanged.  For the
rates, \(N_2\geq nh_2/C_K\) and \(d_2\leq1/2\) give
\(\mathcal E^{(2)}_x\leq3s^{(2)}_1/\lambda_2\) with
\(s^{(2)}_1\leq\sigma L_2\Gamma_1\sqrt{2C_Kz_x/(nh_2)}\).  The bound on
\(\mathcal R^{(1)}_x\) is Note~1, Corollary~4.4, with the step-one
evaluation set equal to the design.
\end{proof}

\end{document}
